\documentclass[12pt,a4paper]{article}

\usepackage[utf8]{inputenc}
\usepackage[T1]{fontenc}
\usepackage[margin=2.5cm]{geometry}
\usepackage{hyperref}
\usepackage{microtype}
\usepackage{parskip}
\usepackage{lmodern}
\usepackage{amsmath}

\hypersetup{colorlinks=true, urlcolor=blue!60!black}

\pagestyle{empty}

\begin{document}

\begin{flushright}
Kundai Farai Sachikonye \\
Auf der Lichtung 19, 82178 Puchheim, Germany \\
\href{https://github.com/fullscreen-triangle}{github.com/fullscreen-triangle} \\
\texttt{kundai.sachikonye@bitspark.com} \\
(+49) 1626386344 \\[6pt]
20 May 2026
\end{flushright}

\vspace{1em}

PD Dr.\ Markus Herrmann \\
Institut für Medizin- und Datenethik (IMDE) \\
Medizinische Fakultät, Universität Heidelberg \\
Tiergartenstraße 15, 69121 Heidelberg \\
Job-ID: V000015525

\vspace{1.5em}

\textbf{Re: Wissenschaftliche*r Mitarbeiter*in, KI- und Datenethik in der Medizin ---
ONCOverse (V000015525)}

\vspace{1em}

Dear PD Dr.\ Herrmann,

I am applying for the research position in AI and data ethics in medicine (V000015525)
in the ONCOverse project.
My background is unusual for an ethics position: I am a computational biologist who
has spent four years building, deploying, and validating AI diagnostic systems for
oncology and immunopeptidology.
I am applying because I believe the gap this project most needs filled is not
between philosophy and medicine, but between the engineers building ONCOverse's
AI models and the ethics team advising them.
That gap requires someone who has sat on both sides.

\textbf{I build the systems being evaluated.}
The ONCOverse AI layer will include diagnostic models for pancreatic cancer
characterisation, prognostic tools that combine multi-omics and clinical data,
and a VR-mediated patient--physician interface generating interaction data.
Ethicists who arrive from philosophy encounter these systems from the outside:
they can critique stated design goals, but the actual ethical risks are buried in
implementation choices that require technical fluency to identify.
I have built directly analogous systems.
The \href{https://syndrome-seven.vercel.app/tools}{Syndrome} framework performs
neoantigen identification, MHC binding prediction (AUC\,$0.81$, validated on 2,847
IEDB peptides), immune escape scoring, and cancer disease state characterisation
via a vector $\mathbf{D}\in[0,1]^8$ with explicit immune activation ($D_\mathrm{Imm}$)
and barrier integrity ($D_\mathrm{Bar}$) dimensions.
The \href{https://github.com/fullscreen-triangle/crown-prince}{Crown Prince} framework
derives pharmacokinetic and pharmacodynamic trajectories from molecular structure alone,
validated at 39/39 against NIST, with therapeutic target identification from
first-principles geometry.
From building these tools I know precisely where the ethical problems arise in
clinical AI: not in the headline claims, but in training data composition,
threshold selection, population generalisability, and the gap between validation
cohort and deployment context.

\textbf{What ``explainability'' actually requires.}
The dominant ethics critique of clinical AI is the black-box problem.
Most ethics guidelines respond with procedural requirements (document the training
data, provide SHAP values) that are technically satisfiable without producing
genuinely auditable systems.
I have built AI tools that are auditable in a stronger sense: the Syndrome
framework uses 3-channel discrete Fourier transforms on amino acid physicochemical
properties, with zero free parameters trained on patient data.
Every prediction is a geometric operation on a known mathematical space,
reproducible from the algorithm alone, independent of the training cohort.
The disease state vector $\mathbf{D}$ is a first-principles normative construct
(not a learned embedding), so its axes carry interpretable meaning that can be
communicated to patients and clinicians.
Translating this experience into ONCOverse ethics guidelines: I can specify what
``explainability'' means technically and write requirements that engineers can
implement rather than satisfy procedurally.

\textbf{Equitable access to ONCOverse innovations.}
The Syndrome toolkit runs in a browser (WebGPU, no server infrastructure required)
and occupies 50\,MB compared to DrugBank's 12\,GB.
This design choice was deliberate: computational diagnostic tools that require
GPU cluster infrastructure are equitable only within institutions that can afford
the infrastructure.
Access justice in the ONCOverse context will require similar analysis ---
not just whether the AI model performs equitably across demographic groups
(algorithmic fairness in the standard sense), but whether the delivery
infrastructure systematically excludes the institutions and patients who
most need it.
I have direct experience quantifying that gap from the design side.

\textbf{Multi-site data governance and patient data ethics.}
The ONCOverse project pools data from numerous university hospitals.
The \href{https://github.com/fullscreen-triangle/gospel}{Gospel} pipeline
I built handles federated multi-omics inference across independent sites:
host genotype, metagenomics, and transcriptomics as a joint Bayesian graph,
validated against 1000~Genomes, gnomAD, and UK~Biobank.
From this I understand concretely what data actually flows in multi-site
clinical AI studies: which fields are re-identifiable, where pseudonymisation
fails, and what the practical constraints on consent architecture are
given the analysis methods being used.
The VR patient networking component of ONCOverse raises additional ethics
questions --- identity disclosure, interaction data retention, therapeutic
context of patient--patient connection --- that are distinct from
standard clinical trial data ethics.
I have analysed analogous security and privacy questions in the context
of distributed AI systems and can bring that analysis to the ONCOverse context.

\textbf{Normative frameworks for clinical AI.}
My research applies formal mathematical structures to biomedical problems,
which requires explicit normative commitments at every modelling decision:
what counts as disease (the $\mathbf{D}$ vector definition),
what counts as a valid therapeutic target (the sparse $\ell_1$ selection criterion),
what precision is sufficient for a diagnostic claim.
These are ethics decisions embedded in computational form.
Writing ethics guidelines for ONCOverse's AI technologies is in part
the inverse problem: extracting the normative commitments that are already
implicit in the technical design, making them explicit, and evaluating
whether they are justifiable.
I am equipped to do that work precisely because I have made those decisions
from the inside.

\textbf{Background.}
I hold an MSc in Computational Biology (Universität Tübingen), an MSc in
Pharmaceutical Biotechnology (HAW Hamburg), and a BSc in Biochemistry and
Cell Biology (Jacobs University Bremen).
I conducted doctoral research in Computational Lipidomics at TUM / Bayerisches
Forschungsinstitut (2019--2021) before transitioning to independent research.
I have lived and worked in Germany continuously since 2011.
English is native; German is conversational.

I am available from 1 July 2026 and welcome the opportunity to discuss
how this background can contribute to the ONCOverse ethics work.

\vspace{1.5em}
Yours sincerely,

\vspace{2.5em}
\textbf{Kundai Farai Sachikonye}

\end{document}
